\section{The two-week market}
\label{sec:dynamics}

Each step advances the state
\(\bigl(p_{t-1},\,q_{i,t-1},\,S_{i,t}\bigr)\) --- last period's world
price, each exporter's offer price, and carry-in stocks --- given that
period's harvest, policy, and demand shifters. The order of operations is
fixed and there is no optimisation nested inside the step:

\begin{enumerate}
\item \textbf{Demand.} Each country forms desired domestic use. Food and
      feed demand is price-responsive with a short-run elasticity; the
      industrial (ethanol) component is not. Demand responds to the
      \emph{previous} period's world price, and there is no separate
      domestic price, so an export ban does not lower consumer prices in
      the restricting country.
\item \textbf{Precautionary cover.} Each country targets stocks sufficient
      to cover consumption until the next harvest, plus a safety margin.
      This is a target-stock rule, not an intertemporal arbitrage
      condition.
\item \textbf{Export supply and import demand.} Surplus above desired use
      and the cover target is offered for export and then reduced by the
      prevailing AMIS restriction. Importers post a food-deficit
      requirement plus partial rebuilding toward their own cover target.
\item \textbf{Bilateral allocation.} Trade follows the observed FAOSTAT
      partner structure, reweighted toward cheaper offer prices, with a
      residual pool when preferred links cannot clear.
\item \textbf{Stock accounting.} Consumption and shipments are settled;
      stocks above physical capacity are drawn down.
\item \textbf{Offer-price adjustment.} Exporters raise offer prices when
      their supply clears, cut them when it does not, and mark up when
      importers' preferred suppliers are restricted.
\item \textbf{World price.} The world price moves toward a target that
      combines the trade-weighted offer price with a scarcity term
      measured against an unshocked counterfactual, and is then smoothed.
\end{enumerate}

Three comparisons are worth stating at the outset, because the sitting
raised each. Agrimate solves per-region constrained programmes at every
step \citep{kuhla2025}; Gate~0 does not. Competitive storage in the
Wright--Williams tradition holds grain when the expected price increase
exceeds the cost of carry \citep{wright1982,williams1991}; Gate~0's cover
rule contains no such arbitrage condition. Spatial price equilibrium
determines price as the multiplier on regional balance
\citep{samuelson1952,takayama1971}; Gate~0 determines it by an explicit
adjustment rule.

The \emph{twin} is a parallel run of the identical mechanism under
climatological harvest, mean food and feed demand, no industrial excess,
and no restrictions. Its price path is not scored. It supplies the
counterfactual against which scarcity is measured, so ``tight'' means
tight relative to a normal year rather than relative to an arbitrary
constant.

\flowfig{01_current_fortnight.pdf}{Order of operations within one
two-week step. Each box is evaluated once; no inner solve.
\label{fig:g0}}

\subsection{Demand and precautionary cover}

Desired domestic use is isoelastic in the lagged world price for the
flexible (food and feed) component and inelastic for the industrial
component:
\begin{equation}
d_{i,t}=C^{\mathrm{flex}}_{i,t}\left(\frac{p_{t-1}}{p_0}\right)^{\varepsilon}
+C^{\mathrm{ind}}_{i,t},
\label{eq:demand}
\end{equation}
where \(d_{i,t}\) is desired demand of country \(i\) in period \(t\)
(MMT per step), \(C^{\mathrm{flex}}_{i,t}\) is food-and-feed demand
evaluated at the reference price (MMT per step), \(C^{\mathrm{ind}}_{i,t}\)
is price-inelastic industrial use (MMT per step; the US maize ethanol
residual, zero elsewhere), \(p_{t-1}\) is the previous period's world
price (2010 \$/t), \(p_0\) is the reference price (2010 \$/t), and
\(\varepsilon<0\) is the short-run own-price elasticity of demand.

Expected harvest blends the current realisation with climatology:
\begin{equation}
H^{\mathrm{exp}}_{i,t}=\phi\,H_{i,t}+(1-\phi)\,H^{\mathrm{seas}}_{i,t},
\label{eq:hexp}
\end{equation}
where \(H^{\mathrm{exp}}_{i,t}\) is expected harvest (MMT per step),
\(H_{i,t}\) is realised harvest (MMT per step), \(H^{\mathrm{seas}}_{i,t}\)
is the climatological harvest path (MMT per step), and
\(\phi\in[0,1]\) is the weight on current information. Expectations are
adaptive in this single sense; there is no multi-year forecast.

The stock target is lean-season cover plus a safety margin:
\begin{equation}
T_{i,t}=L_{i,t}+s_i,
\qquad
s_i=\sigma^{\mathrm{stu}}\,C^{\mathrm{ann}}_i,
\label{eq:target}
\end{equation}
where \(T_{i,t}\) is target stock (MMT), \(L_{i,t}\) is the shortfall of
expected harvest against consumption over the interval to the next world
harvest pulse (MMT), \(s_i\) is the safety stock (MMT),
\(\sigma^{\mathrm{stu}}\) is the target stock-to-use ratio (years of use),
and \(C^{\mathrm{ann}}_i\) is annual domestic use (MMT per year). The
horizon is defined by the arrival of a fixed fraction of the annual
\emph{world} harvest, so it is common across countries.

\subsection{Export supply, import demand, and allocation}

Export supply is residual surplus, net of policy:
\begin{equation}
O_{i,t}=\max\bigl(0,\;\mathrm{avail}_{i,t}-d_{i,t}-T_{i,t}\bigr)\,(1-\tau_{i,t}),
\label{eq:offers}
\end{equation}
where \(O_{i,t}\) is export supply offered (MMT per step),
\(\mathrm{avail}_{i,t}=S_{i,t}+H_{i,t}\) is grain available (MMT),
\(S_{i,t}\) is carry-in stock (MMT), and \(\tau_{i,t}\in[0,1]\) is the
export restriction intensity from the AMIS policy record (a ban maps to
\(\approx 0.95\), an export tax to \(\approx 0.50\)).

Import demand combines an immediate deficit with partial rebuilding:
\begin{equation}
D_{i,t}=\underbrace{\max\bigl(0,\,d_{i,t}-\mathrm{avail}_{i,t}\bigr)}_{\text{consumption deficit}}
+\lambda\,\max\bigl(0,\,T_{i,t}-\max(0,\mathrm{avail}_{i,t}-d_{i,t})\bigr),
\label{eq:importdemand}
\end{equation}
where \(D_{i,t}\) is import demand (MMT per step) and \(\lambda\) is the
partial-adjustment rate toward the stock target (per step). Equation
\eqref{eq:importdemand} is a stock-adjustment specification, not a
solution to a dynamic programme.

Bilateral allocation is Armington on a prescribed network
\citep{armington1969}. Preferred destination shares are reweighted toward
low offer prices,
\begin{equation}
\tilde A_{ij,t}\;\propto\;A_{ij}\left(\frac{p_0}{q_{i,t-1}}\right)^{\gamma},
\qquad \sum_j \tilde A_{ij,t}=1,
\label{eq:reweight}
\end{equation}
where \(A_{ij}\) is the observed FAOSTAT destination share of exporter
\(i\) to importer \(j\) (dimensionless), \(q_{i,t-1}\) is exporter \(i\)'s
lagged offer price (2010 \$/t), and \(\gamma>0\) governs how strongly
buyers shift toward cheaper origins. Shipments on each link are the short
side,
\begin{equation}
X_{ij,t}=\min\bigl(O_{i,t}\tilde A_{ij,t},\;S_{ij}D_{j,t}\bigr),
\label{eq:clear}
\end{equation}
where \(X_{ij,t}\) is the shipment from \(i\) to \(j\) (MMT per step) and
\(S_{ij}\) is the observed source share of importer \(j\) from exporter
\(i\) (dimensionless; note \(S_{ij}\) is a share, whereas \(S_{i,t}\) is a
stock). Unmatched supply then meets a fraction \(\nu\) of unmatched demand
through a residual pool. Quantities therefore come from
\eqref{eq:offers}--\eqref{eq:importdemand}; the trade matrix supplies only
the pattern of partners.

\subsection{Stock accounting and capacity}

Realised consumption is bounded by post-trade availability:
\begin{equation}
C^{\mathrm{used}}_{i,t}=\min\Bigl(d_{i,t},\;
\max\bigl(0,\,\mathrm{avail}_{i,t}-X^{\mathrm{out}}_{i,t}+X^{\mathrm{in}}_{i,t}\bigr)\Bigr),
\label{eq:consumption}
\end{equation}
where \(C^{\mathrm{used}}_{i,t}\) is realised consumption (MMT per step),
\(X^{\mathrm{out}}_{i,t}=\sum_j X_{ij,t}\) is total exports and
\(X^{\mathrm{in}}_{i,t}=\sum_j X_{ji,t}\) total imports (MMT per step).
Closing stock is the residual, less any drawdown of grain held above
physical capacity:
\begin{equation}
S_{i,t+1}=S^{-}_{i,t+1}-\lambda_W\max\bigl(0,\,S^{-}_{i,t+1}-W_{i,t}\bigr),
\qquad
S^{-}_{i,t+1}=\max\bigl(0,\,\mathrm{avail}_{i,t}-X^{\mathrm{out}}_{i,t}
+X^{\mathrm{in}}_{i,t}-C^{\mathrm{used}}_{i,t}\bigr),
\label{eq:stocks}
\end{equation}
where \(S_{i,t+1}\) is closing stock (MMT), \(S^{-}_{i,t+1}\) is closing
stock before the capacity adjustment (MMT), \(W_{i,t}\) is storage
capacity (MMT), and \(\lambda_W\) is the rate at which excess above
capacity is drawn down (per step). This term removes grain physically; it
is not a pecuniary cost of carry, and no interest rate appears anywhere in
\eqref{eq:stocks}. Capacity is
\begin{equation}
W_{i,t}=\sigma^{\max}C^{\mathrm{ann}}_i
+\frac{n^{\mathrm{pipe}}}{24}\,C^{\mathrm{ann}}_i
+\mathbf{1}\{n^{\mathrm{pipe}}>0\}\,H_{i,t},
\label{eq:capacity}
\end{equation}
where \(\sigma^{\max}\) is the maximum stock-to-use ratio (years of use)
and \(n^{\mathrm{pipe}}\) is working inventory in steps (12 for wheat and
maize, i.e.\ six months; 0 for rice, which harvests year-round).

\subsection{Offer prices}

Exporters adjust offer prices toward market conditions:
\begin{equation}
q_{i,t}=(1-\beta)\,q_{i,t-1}
\exp\Bigl(\alpha\bigl(f_{i,t}-\theta\bigr)+\alpha_r\,b_t\,\mathbf{1}\{O_{i,t}>0\}\Bigr)
+\beta\,p_{t-1},
\label{eq:ask}
\end{equation}
where \(q_{i,t}\) is exporter \(i\)'s offer price (2010 \$/t),
\(f_{i,t}=X^{\mathrm{out}}_{i,t}/O_{i,t}\) is the fill rate
(dimensionless), \(\theta\) is the fill rate at which the price is left
unchanged, \(\alpha>0\) is the adjustment gain (per step), \(b_t\) is the
share of world import demand whose preferred suppliers are restricted
(dimensionless), \(\alpha_r\ge 0\) is the markup unrestricted exporters
apply when rivals are blocked, and \(\beta\in[0,1]\) is the weight pulling
offer prices toward the world price. Offer prices are bounded to
\([0.45p_0,\,2.8p_0]\).

Equation \eqref{eq:ask} is an adjustment rule: excess demand raises the
price, excess supply lowers it. It is not the first-order condition of an
exporter's profit-maximisation problem, and \(\alpha,\theta,\beta\) are
not structural markups.

\subsection{The world price}

Define market-accessible stocks --- total stocks net of lean-season
requirements and of grain sequestered behind restrictions:
\begin{equation}
\mathcal{F}_t=\sum_i S_{i,t+1}-\sum_i L_{i,t}
-\sum_i \tau_{i,t}\max\bigl(0,\,S_{i,t+1}-T_{i,t}\bigr),
\label{eq:free}
\end{equation}
where \(\mathcal{F}_t\) is market-accessible stock (MMT). The third term
matters: without it, a restriction would register as abundance. Let
\(\mathcal{F}^{\mathrm{twin}}_t\) be the same quantity in the twin, and
define the scarcity ratio
\begin{equation}
r_t=\frac{\mathcal{F}^{\mathrm{twin}}_t+f}{\mathcal{F}_t+f},
\label{eq:ratio}
\end{equation}
where \(f>0\) is a small regularising constant (MMT) that keeps
\eqref{eq:ratio} finite when accessible stocks approach zero. The price
target and its updating are
\begin{align}
p^{\mathrm{scar}}_t&=p_0\,r_t^{\eta}\bigl(1+\kappa_u\Delta u_t+\kappa_b b_t\bigr),
\label{eq:pscar}\\[2pt]
p^{\star}_t&=\omega\,p^{\mathrm{tr}}_t+(1-\omega)\,p^{\mathrm{scar}}_t,
\label{eq:pstar}\\[2pt]
p_t&=\rho\,p_{t-1}+(1-\rho)\,p^{\star}_t,
\label{eq:pupdate}
\end{align}
where \(p^{\mathrm{scar}}_t\) is the scarcity price (2010 \$/t),
\(\eta>0\) is the elasticity of price with respect to the inverse of
accessible stocks, \(\Delta u_t\) is the excess of unmet import demand over
its twin value (dimensionless), \(\kappa_u,\kappa_b\ge 0\) weight unmet
demand and restriction blockage, \(p^{\mathrm{tr}}_t\) is the
shipment-weighted mean offer price (2010 \$/t), \(\omega\in[0,1]\) is the
weight on realised trade prices, \(p^{\star}_t\) is the within-period price
target (2010 \$/t), and \(\rho\in[0,1]\) is the smoothing parameter. The
world price is bounded to \([60,1200]\) \$/t.

Equation \eqref{eq:pscar} is inverse-demand in form: price rises as
accessible stocks fall relative to normal, with elasticity \(\eta\). The
exponent is symmetric, so a surplus lowers the price and food and feed
demand absorbs it; muting the abundance side left gluts sitting in store.
Equations \eqref{eq:pstar}--\eqref{eq:pupdate} then mix that term with
realised trade prices and smooth the result, so the world price is not
required to lie on the inverse-demand schedule within the period. That is
the honest gap, and it is the subject of \S\ref{sec:options}.

\subsection{Exact properties}

\paragraph{Reference identity.}
If the crisis run matches the twin in accessible stocks, unmet import
demand, and blockage --- \(\mathcal{F}_t=\mathcal{F}^{\mathrm{twin}}_t\),
\(\Delta u_t=0\), \(b_t=0\) --- then \(r_t=1\) and \eqref{eq:pscar}
gives \(p^{\mathrm{scar}}_t=p_0\). With offer prices at \(p_0\),
\eqref{eq:pstar} gives \(p^{\star}_t=p_0\) exactly, and
\eqref{eq:pupdate} is a geometric approach to \(p_0\). This is an algebraic
identity, not a calibrated residual. Note that \(\mathcal{F}^{\mathrm{twin}}_t\)
itself oscillates seasonally: the reference is a trajectory, not a constant.

\paragraph{Twin closure.}
During the twin run the price is held at \(p_0\), so
\(\mathcal{F}^{\mathrm{twin}}_t\) reflects physical availability and does
not feed back through \eqref{eq:demand}. The twin is a control, not a
second priced market.

\paragraph{Mass balance.}
Availability is exhausted by consumption, net trade, the change in stocks,
and the capacity drawdown in \eqref{eq:stocks}. The drawdown is part of the
specification and is reported, not absorbed silently.
